\chapter{Pipeline 五路径路由}\label{ch:routing}

\section{路由总决策树}

\texttt{RunBackupPipeline}（\filepath{backup\_pipeline.cc}）在备份开始时按文件数、大小、模式、worker 配置选择拓扑。

\begin{figure}[htbp]
\centering
\begin{tikzpicture}[node distance=0.55cm, scale=0.88, transform shape]
  \node[block=CoreBlue, minimum width=3.8cm] (entry) {RunBackupPipeline};
  \node[decision, below=of entry] (n1) {文件数=1?};
  \node[decision, below left=1.0cm and 1.5cm of n1] (multi) {多文件 / 增量 /\\workers$>$0?};
  \node[decision, below right=1.0cm and 1.5cm of n1] (single) {单文件};
  \node[decision, below=of single] (sz) {total $\leq$ 32MB?};
  \node[block=PathD, below left=0.9cm and 0.3cm of sz, minimum width=3.2cm] (inline) {L3a Inline\\RunSingleFileInline};
  \node[decision, below right=0.9cm and 0.3cm of sz] (big) {$>$32MB 且\\Full 且 workers=0?};
  \node[block=PathA, below=of big, minimum width=3.5cm] (stream) {Streaming\\RunSingleFileStreaming};
  \node[block=PathB, below=of multi, minimum width=3.5cm] (v4) {Pipeline v4\\多 worker 全局队列};
  \draw[arr] (entry) -- (n1);
  \draw[arr] (n1) -- node[lab,left]{否} (multi);
  \draw[arr] (n1) -- node[lab,right]{是} (single);
  \draw[arr] (single) -- (sz);
  \draw[arr] (sz) -- node[lab,left]{是} (inline);
  \draw[arr] (sz) -- node[lab,right]{否} (big);
  \draw[arr] (big) -- node[lab,right]{是} (stream);
  \draw[arr] (big.east) -| ++(1.5,0) |- node[lab,right,near start]{否} (v4.east);
  \draw[arr] (multi) -- (v4);
\end{tikzpicture}
\caption{顶层路由决策树（\filepath{backup\_pipeline.cc:1393--1424}）。}
\label{fig:top-router}
\end{figure}

\section{五路径详解}

\begin{table}[htbp]
\centering
\caption{五路径对比一览}
\label{tab:five-paths}
\footnotesize
\begin{tabular}{@{}llllp{3.8cm}@{}}
\toprule
路径 & 触发条件 & CDC 引擎 & Worker & 典型场景 \\
\midrule
\textbf{L3a Inline} & 1 文件 $\leq$32MB & 内联 FastCDC & 0（单线程） & 小文件快路径 \\
\textbf{Streaming} & 1 文件 $>$32MB Full & StreamFeed 32MB & 2 enc + 1 store & L5 大文件 \\
\textbf{Hybrid} & Streaming 内默认 & StreamFeed（同左） & 同左 & Wave I 默认 \\
\textbf{FastSlice} & \texttt{CDC\_FAST\_SLICE=1} & 整文件 ChunkCuts & 同 Streaming 拓扑 & opt-in，慢 \\
\textbf{Pipeline v4} & 多文件/增量/workers & HCRBO 或 Stream & N$\times$R/C/E/S & L7 multi-file \\
\bottomrule
\end{tabular}
\end{table}

\section{Streaming 内部子路由}

单文件进入 \texttt{RunSingleFileStreamingChunkPipeline} 后，\texttt{ChunkFileStreaming} 再次分支：

\begin{figure}[htbp]
\centering
\begin{tikzpicture}[node distance=0.55cm]
  \node[block=CoreBlue] (cs) {ChunkFileStreaming};
  \node[decision, below=of cs] (hyb) {Hybrid 开启?\\（默认是）};
  \node[block=PathA, below left=0.9cm and 0.6cm of hyb] (hpath) {ChunkFileStreamingHybrid};
  \node[decision, below right=0.9cm and 0.6cm of hyb] (fast) {FastSlice opt-in?};
  \node[block=PathC, below=of fast] (fpath) {ChunkFileStreamingFastSlice};
  \node[block=PathA, below left=0.9cm and 0.6cm of fast] (spath) {ChunkFileStreamingFeed};
  \draw[arr] (cs) -- (hyb);
  \draw[arr] (hyb) -- node[lab,left]{是} (hpath);
  \draw[arr] (hyb) -- node[lab,right]{否} (fast);
  \draw[arr] (fast) -- node[lab,right]{是} (fpath);
  \draw[arr] (fast) -- node[lab,left]{否} (spath);
\end{tikzpicture}
\caption{Streaming 内部 CDC 子路由（\filepath{backup\_pipeline.cc:799--813}）。}
\label{fig:stream-sub-router}
\end{figure}

路由优先级：\textbf{Hybrid $\rightarrow$ FastSlice $\rightarrow$ Stream Feed}。

\section{各路径线程模型}

\subsection{L3a Inline（\texorpdfstring{$\leq$32MB}{<=32MB}）}

\begin{figure}[htbp]
\centering
\begin{tikzpicture}[node distance=0.4cm]
  \node[block=PathD, minimum width=8cm] (t1) {单线程：Read $\rightarrow$ Chunk $\rightarrow$ Encode $\rightarrow$ Store};
\end{tikzpicture}
\caption{无队列、无 worker；开销最低。}
\end{figure}

\subsection{StreamingChunkCpuPipeline（\texorpdfstring{$>$32MB Full}{$>$32MB Full}）}

\begin{figure}[htbp]
\centering
\begin{tikzpicture}[node distance=0.5cm]
  \node[block=PathA, minimum width=3.2cm] (main) {主线程\\mmap + StreamFeed\\CDC + digest};
  \node[block=CoreOrange, right=1.2cm of main, minimum width=2.2cm] (enc) {Encode$\times$2};
  \node[block=CoreGreen, right=1.0cm of enc, minimum width=2.0cm] (sto) {Store$\times$1};
  \draw[arr] (main) -- node[lab,above]{ChunkTask Q} (enc);
  \draw[arr] (enc) -- node[lab,above]{Encoded Q} (sto);
\end{tikzpicture}
\caption{固定 2+1 worker；队列深度 1024。}
\end{figure}

\subsection{Pipeline v4（多文件 / 增量）}

\begin{figure}[htbp]
\centering
\begin{tikzpicture}[node distance=0.45cm, scale=0.9, transform shape]
  \node[block=PathB, minimum width=2.5cm] (sched) {FileScheduler};
  \node[block=CoreBlue, below left=0.7cm and 0.5cm of sched] (r1) {Reader$\times$N};
  \node[block=CoreOrange, below=0.7cm of sched] (c1) {Chunker$\times$N};
  \node[block=CorePurple, below right=0.7cm and 0.5cm of sched] (e1) {Encode$\times$M};
  \node[block=CoreGreen, below=1.5cm of c1] (s1) {Store$\times$K};
  \node[block=CoreBlue, below=0.5cm of s1] (agg) {FileAggregator};
  \draw[arr] (sched) -- (r1);
  \draw[arr] (r1) -- (c1);
  \draw[arr] (c1) -- (e1);
  \draw[arr] (e1) -- (s1);
  \draw[arr] (s1) -- (agg);
\end{tikzpicture}
\caption{v4 全局有界队列；多文件并行；增量走 HCRBO。}
\end{figure}

\section{为何增量必须走 v4？}

\begin{corebox}[路由分化理由]
单文件 Streaming 路径专为 \textbf{full backup 大文件} 的 feed overlap 优化。\\
增量备份需要：
\begin{itemize}[nosep]
  \item \texttt{EbHcrboChunker::ChunkIncremental} + prior CFI；
  \item 多文件 \texttt{FileScheduler} 调度；
  \item 与 Pipeline v4 的 \texttt{ChunkTask} 队列/manifest 汇总集成。
\end{itemize}
故增量\textbf{不}走 \texttt{RunSingleFileStreamingChunkPipeline}。
\end{corebox}

\section{环境变量速查}

\begin{table}[htbp]
\centering
\caption{CDC / Pipeline 路由相关环境变量}
\label{tab:env}
\small
\begin{tabular}{@{}llp{6.5cm}@{}}
\toprule
变量 & 默认 & 效果 \\
\midrule
\texttt{EBBACKUP\_CDC\_HYBRID} & 开 & Hybrid stream feed \\
\texttt{EBBACKUP\_CDC\_FAST\_SLICE} & 关 & Phase A 整文件 cuts \\
\texttt{EBBACKUP\_FORCE\_STREAM\_CDC} & 关 & 禁 FastSlice，强制 Feed \\
\texttt{EBBACKUP\_PIPELINE\_WORKERS} & — & 显式 worker $\Rightarrow$ 走 v4 \\
\texttt{EBBACKUP\_PIPELINE\_PROFILE} & 关 & 打印 phase 毫秒分解 \\
\bottomrule
\end{tabular}
\end{table}

\section{性能画像：为何路由分化而非 Pipeline v5}

\begin{figure}[htbp]
\centering
\begin{tikzpicture}
\begin{axis}[
  width=0.88\linewidth,
  height=6cm,
  ybar,
  bar width=18pt,
  ymin=0, ymax=700,
  ylabel={吞吐 (MB/s)},
  symbolic x coords={L3a,L5 Stream,L5 FastSlice,L7 v4},
  xtick=data,
  nodes near coords,
  every node near coord/.append style={font=\scriptsize},
  grid=major,
  grid style={dashed,gray!30},
  legend style={at={(0.5,-0.22)},anchor=north,font=\footnotesize}
]
\addplot[fill=PathD!50,draw=PathD] coordinates {(L3a,120) (L5 Stream,171) (L5 FastSlice,105) (L7 v4,650)};
\addlegendentry{典型 bench 量级（256MB/32$\times$32MB）}
\end{axis}
\end{tikzpicture}
\caption{不同路径服务不同 bench 维度；分化而非统一 Pipeline v5。}
\label{fig:bench-paths}
\end{figure}

\begin{remark}
v0.8 Pipeline v4 把 L7 拉到 550+ MB/s，但 L5 单文件仍 $\sim$146 MB/s（chunk 占 98\%）。\\
v0.9.2 起\textbf{按场景路由}到 Streaming，四轮 sprint 把 L5 拉到 171 MB/s，且 L7 几乎不变。
\end{remark}
